% =============================================================================
% Chapter 4 --- RTL Module Design
% Grounded in Context 1 §3-5 (modules deep dive) and Context 2 entirely
% (FSM details). All signal names and transition conditions come from the
% SystemVerilog source files listed in those Context documents.
% =============================================================================

\chapter{RTL Module Design}
\label{ch:rtl}

This chapter is the technical heart of the report. It describes, at the
register-transfer level, the three synthesizable modules of the controller
stack. For each module the discussion follows the actual signal flow in
the source file: continuous assignments, registered captures,
combinational next-state logic, and output multiplexing. FSM state
diagrams in this chapter are deliberately simplified for readability;
fully-labeled versions are collected in \Cref{app:fsm-diagrams}.

\section{Parallel Interface}
\label{sec:rtl-pi}

\subsection{Module Structure}
\label{sec:rtl-pi-structure}

The \modname{parallel\_interface} module is described in
\sv{source/Parallel\_Interface/parallel\_interface.sv}. Its body contains
only continuous assignments: the clock and reset ports are declared but
are never referenced inside the module, so no sequential logic is
inferred. The four output assignments are:
\begin{align*}
    \signame{data}    &\leftarrow \signame{host\_data}[31:24] \\
    \signame{address} &\leftarrow \texttt{ann\_tail\_to\_parallel\_addr}(\signame{host\_data}[23:0]) \\
    \signame{cmd}     &\leftarrow \signame{host\_cmd} \\
    \signame{valid}   &\leftarrow (\signame{host\_cmd} \neq \sv{CMD\_HIZ})
                                   \land \texttt{ann\_tail\_is\_valid\_onehot}(\signame{host\_data}[23:0])
\end{align*}
Any synchronization of host-side signals must therefore be performed
outside this module; the RTL in scope is purely combinational.

\subsection{One-Hot Tail Validation}
\label{sec:rtl-pi-onehot}

The \signame{valid} qualifier is the logical AND of (i)~a non-idle command
and (ii)~a strictly one-hot tail. Tail legality is the AND of four
independent structural predicates, as summarized in
\Cref{tab:pi-onehot-predicates}.

\begin{table}[htbp]
    \centering
    \caption{One-hot predicates used to form \signame{valid} in
        \modname{parallel\_interface}.}
    \label{tab:pi-onehot-predicates}
    \begin{tabularx}{\linewidth}{l c l X}
        \toprule
        \textbf{Tail slice} & \textbf{Width} & \textbf{Predicate} & \textbf{Meaning} \\
        \midrule
        \sv{tail[23:20]} & 4 & \sv{pi\_is\_onehot4} & Exactly one PE-select bit \\
        \sv{tail[19:16]} & 4 & \sv{pi\_is\_onehot4} & Exactly one SA-select bit \\
        \sv{tail[15:8]}  & 8 & \sv{pi\_is\_onehot8} & Exactly one column bit \\
        \sv{tail[7:0]}   & 8 & \sv{pi\_is\_onehot8} & Exactly one row bit \\
        \bottomrule
    \end{tabularx}
\end{table}

\sv{pi\_is\_onehot4} returns 1 only for the four patterns
\sv{4'b0001}, \sv{4'b0010}, \sv{4'b0100}, \sv{4'b1000}; any other value
(including all-zeros or multi-hot) returns 0. \sv{pi\_is\_onehot8}
similarly returns 1 only for the eight single-bit patterns and 0
otherwise. Consequently, a host packet with a non-\sv{HIZ} command but
any non-one-hot tail holds \signame{valid} low even though \signame{data}
and \signame{address} are still driven combinationally from
\signame{host\_data}. The downstream controller gates acceptance of new
transactions on \signame{valid}, so invalid packets are silently
rejected at the interface.

\subsection{Tail-to-Address Mapping}
\label{sec:rtl-pi-addr}

The function \sv{ann\_tail\_to\_parallel\_addr} performs a
\emph{decode-then-concatenate} path entirely inside the package. First,
four case-style decoders produce packed indices:

\begin{align*}
    \texttt{blk} &= \texttt{pi\_onehot4\_to\_idx}(\sv{tail[23:20]}) \in [0,3] \\
    \texttt{sb}  &= \texttt{pi\_onehot4\_to\_idx}(\sv{tail[19:16]}) \in [0,3] \\
    \texttt{cid} &= \texttt{pi\_onehot8\_to\_idx}(\sv{tail[15:8]})  \in [0,7] \\
    \texttt{rid} &= \texttt{pi\_onehot8\_to\_idx}(\sv{tail[7:0]})   \in [0,7]
\end{align*}

\noindent Then these indices are concatenated into a 16-bit address:
\[
    \signame{address} = \{\sv{6'b0},~\texttt{blk},~\texttt{sb},~\texttt{cid},~\texttt{rid}\}
\]

The returned 16-bit value therefore carries \sv{[9:8]}~=~\texttt{blk},
\sv{[7:6]}~=~\texttt{sb}, \sv{[5:3]}~=~\texttt{cid},
\sv{[2:0]}~=~\texttt{rid}, with the upper six bits always zero. This is
the exact field layout that \sv{parse\_ann\_address} later expects inside
the controller (\Cref{sec:rtl-controller-addr}). The inverse operation,
indices to one-hot tail, is implemented in \sv{build\_host\_ann\_word}
via \sv{4'(1 << blk)}, \sv{8'(1 << cid)}, and so on, giving a structural
dual of the decode path.

\subsection{Relation to the Controller}
\label{sec:rtl-pi-controller}

The controller does not re-validate one-hot: it consumes \signame{address}
as a plain bit-field slice. Semantic consistency between the host's
one-hot tail and the resulting \signame{address} is enforced at the PI
boundary when \signame{valid} is high, not inside the FSM. Invalid tails
never cause the controller to enter a new state.

\section{ANN Controller}
\label{sec:rtl-controller}

The controller is declared as \modname{ann\_controller} in
\sv{source/Controller/controller.sv}. It is parameterized by
\sv{ADDR\_WIDTH}, \sv{WEIGHT\_WIDTH}, \sv{USE\_WEIGHT\_PULSE\_LUT}, and
\sv{WEIGHT\_PULSE\_LUT\_FILE}. The module contains one large combinational
block (wrapped by the \sv{\textasciigrave comb} macro from
\sv{source/common/macros.svh}) that assigns all outputs and all
\signame{next\_*} values, together with two \sv{always\_ff} blocks that
register the main FSM state, the sub-FSM states, counters, and captured
address/data.

\subsection{Registered Capture of Host Data}
\label{sec:rtl-controller-capture}

The controller latches the host bus only at the transition out of
\stateval{S\_IDLE}. In its sequential block, on the positive edge of
\signame{clk} with asynchronous active-low reset:

\begin{itemize}
    \item if \signame{valid} is high and the main state is
          \stateval{S\_IDLE}, the RTL performs
          \signame{address\_reg} \(\leftarrow\) \signame{address} and
          \signame{data\_reg} \(\leftarrow\) \signame{data};
    \item otherwise the registers retain their values.
\end{itemize}

Once captured, all subsequent sub-FSM states (program, verify, read,
erase) use the \emph{registered} \signame{address\_reg} and
\signame{data\_reg} rather than the live PI outputs. The only exception
is the inference path, where the combinational FSM explicitly uses the
live \signame{data} byte for pixel streaming during
\stateval{S\_COLLECT\_DATA} (see \Cref{sec:rtl-controller-mainfsm}).

\subsection{Address Parsing and Word Packing}
\label{sec:rtl-controller-addr}

A combinational region calls the package function
\sv{parse\_ann\_address} on \signame{address\_reg} and deposits the
resulting indices into a 10-bit local \signame{weight\_addr\_reg}:

\begin{align*}
    \texttt{block\_id}     &= \signame{address\_reg}[9:8] \\
    \texttt{sub\_block\_id} &= \signame{address\_reg}[7:6] \\
    \texttt{col\_id}       &= \signame{address\_reg}[5:3] \\
    \texttt{row\_id}       &= \signame{address\_reg}[2:0]
\end{align*}

Bits \sv{[15:10]} of the address are not consumed. When the state is
not \stateval{S\_IDLE}, the outbound \signame{ann\_core\_word} is produced
combinationally by
\sv{pack\_ann\_core\_word(data\_byte\_for\_ann, block\_id, sub\_block\_id,
row\_id, col\_id)}, which internally calls
\sv{host\_addr\_to\_ann\_addr\_out} to build each one-hot field by a
left-shift:

\begin{align*}
    \texttt{pe\_onehot}  &= 1 \ll \texttt{block\_id} \\
    \texttt{sa\_onehot}  &= 1 \ll \texttt{sub\_block\_id} \\
    \texttt{col\_onehot} &= 1 \ll \texttt{col\_id} \\
    \texttt{row\_onehot} &= 1 \ll \texttt{row\_id}
\end{align*}

The final 32-bit word is
\sv{\{data\_byte, pe\_onehot, sa\_onehot, col\_onehot, row\_onehot\}}.
No ``ANN tail register'' is stored; the outbound word is recomputed every
cycle from the current \signame{address\_reg} parse and the selected
payload byte. When the state is \stateval{S\_IDLE}, the RTL forces
\signame{ann\_core\_word} to all zeros.

\subsection{Payload Byte Multiplexing}
\label{sec:rtl-controller-data}

The byte inserted into the outbound word,
\signame{data\_byte\_for\_ann}, is selected in an \sv{always\_comb} with
\sv{unique case (state)}:

\begin{itemize}
    \item \stateval{S\_PROGRAM}, \stateval{S\_VERIFY}, \stateval{S\_READ},
          \stateval{S\_ERASE}: use the captured \signame{data\_reg}.
    \item \stateval{S\_COLLECT\_DATA}: use the live \signame{data}
          coming from the parallel interface, so that streamed
          inference pixels are forwarded to the core.
    \item Default: \sv{8'b0}.
\end{itemize}

\subsection{Main FSM}
\label{sec:rtl-controller-mainfsm}

The main state machine has nine states. The combinational block sets the
following defaults before the \sv{unique case (state)}: \signame{ann\_reset}
and \signame{busy} default to 0; \signame{buf\_reg\_add} defaults to the
(zero-extended) \signame{buf\_addr\_reg}; \signame{buf\_reg\_ctrl} defaults
to \sv{CTRL\_IDLE}; \signame{buf\_read\_write} defaults to read;
\signame{buf\_bit\_sel} is \signame{bit\_count} when the state is
\stateval{S\_COMPUTE} and \sv{3'd0} otherwise; \signame{next\_state} is
\signame{state}; and \signame{next\_prog\_state} is \signame{prog\_state}.
Per-state arms override these defaults where needed.

\Cref{fig:main-fsm-body} summarizes the allowed transitions; a fully
labeled version appears in \Cref{fig:app-main-fsm-full}.

\begin{figure}[htbp]
    \centering
    \includegraphics[width=0.95\linewidth]{fig_main_fsm_body}
    \caption[Main FSM of the ANN controller (simplified)]{Simplified main
        FSM of \modname{ann\_controller}. Arrows show the condition under
        which \signame{next\_state} is set. The labels on transitions out
        of \stateval{S\_VERIFY} and \stateval{S\_ERASE} summarize the
        compound conditions from the respective sub-FSMs.}
    \label{fig:main-fsm-body}
\end{figure}

\paragraph{\stateval{S\_IDLE}.}
\signame{busy} is low; the buffer interface is set to
\sv{CTRL\_IDLE} with \signame{buf\_reg\_add} forced to zero. If
\signame{valid} is low, \signame{next\_state} stays \stateval{S\_IDLE}.
If \signame{valid} is high, \sv{unique case (cmd)} decodes the host
command: \sv{CMD\_PROG} sets \signame{next\_state} to \stateval{S\_PROGRAM}
and \signame{next\_prog\_state} to \sv{PROG\_HIZ}; \sv{CMD\_ERASE} goes
to \stateval{S\_ERASE}; \sv{CMD\_READ} goes to \stateval{S\_READ};
\sv{CMD\_INF} goes to \stateval{S\_COLLECT\_DATA}; the default keeps
\stateval{S\_IDLE}. Only the PROG arm explicitly updates
\signame{next\_prog\_state} within the idle branch; the other arms leave
it at its default.

\paragraph{\stateval{S\_RESET}.}
\signame{busy} and \signame{ann\_reset} are asserted; the buffer is
held in \sv{CTRL\_IDLE}. The transition out is unconditional:
\signame{next\_state} becomes \stateval{S\_PROGRAM} with
\signame{next\_prog\_state}~=~\sv{PROG\_HIZ}. No branch of the
main FSM currently drives \signame{next\_state} into \stateval{S\_RESET};
the state is present in the RTL and drives \signame{ann\_reset} for one
cycle if entered, but in the published FSM graph it is not reachable from
\stateval{S\_IDLE}. Legacy initialization code in the sequential block
remains guarded by
\signame{state}~=~\stateval{S\_IDLE}~\sv{\&\&}~\signame{next\_state}~=~\stateval{S\_RESET},
and is therefore dead in the current graph; this is discussed further in
\Cref{ch:limitations}.

\paragraph{\stateval{S\_PROGRAM}.}
\signame{busy} is high. Buffer control is multiplexed by
\signame{prog\_state}: during \sv{PROG\_HIZ} and \sv{PROG\_SELECT} the
controller writes the data byte into the buffer at
\signame{address\_reg}[5:0] under \sv{CTRL\_DATA\_LOAD}; in all other
sub-states (\sv{PROG\_WRITE}, \sv{PROG\_WAIT\_ACK}, \sv{PROG\_COMPLETE},
and the default) the buffer is configured for \sv{CTRL\_WEIGHT\_READ} at
the same address. The main-state transition is driven by the PROG
sub-FSM: when \signame{prog\_state}~=~\sv{PROG\_COMPLETE}, the shared
combinational block sets \signame{next\_state}~=~\stateval{S\_VERIFY} and
\signame{next\_prog\_state}~=~\sv{PROG\_HIZ}.

\paragraph{\stateval{S\_VERIFY}.}
\signame{busy} is high; the buffer is read under \sv{CTRL\_WEIGHT\_READ}
at \signame{address\_reg}[5:0]. Transition rules are implemented in the
VERIFY sub-FSM (\Cref{sec:rtl-controller-verify}): on
\sv{VERIFY\_DONE}, \signame{next\_state} becomes \stateval{S\_IDLE}; on
\sv{VERIFY\_CHECK} with a readback that is less than expected and
\signame{prog\_retry\_cnt}~\(<\)~\sv{MAX\_PROG\_RETRIES}, the main state
goes back to \stateval{S\_PROGRAM}; on the over-programmed branch or
when retries are exhausted, it goes to \stateval{S\_ERASE}.

\paragraph{\stateval{S\_ERASE}.}
\signame{busy} is high; the buffer is held in \sv{CTRL\_IDLE}. The ERASE
sub-FSM (\Cref{sec:rtl-controller-erase}) drives exit either to
\stateval{S\_IDLE} or back to \stateval{S\_PROGRAM}, depending on the
\signame{erase\_from\_host} flag and the \signame{retry\_cnt} value.

\paragraph{\stateval{S\_READ}.}
\signame{busy} is high; the buffer is held in \sv{CTRL\_IDLE}. The
state self-loops while \signame{pulse\_done} is low, and returns to
\stateval{S\_IDLE} when \signame{pulse\_done} asserts. The
\signame{pulses} bus carries \sv{PULSE\_MODE\_READ} while
\sv{pulse\_train\_active} is true for \stateval{S\_READ}.

\paragraph{\stateval{S\_COLLECT\_DATA}.}
\signame{busy} is high. Buffer control is
\sv{CTRL\_DATA\_LOAD} with \signame{buf\_read\_write}~=~1 and
\signame{buf\_reg\_add}~=~\signame{row\_count}\(\times\)8\(+\)\signame{data\_count}.
Transitions depend on \signame{valid}:
if \signame{valid} is high and \signame{data\_count}~\(<\)~7, the state
self-loops; if \signame{valid} is high and \signame{data\_count}~=~7, the
next state is \stateval{S\_COMPUTE}; if \signame{valid} is low, the state
self-loops. In the sequential block, the relevant counters
(\signame{data\_count}, \signame{row\_count}, \signame{buf\_write\_addr})
are updated only on cycles with \signame{valid} asserted.

\paragraph{\stateval{S\_COMPUTE}.}
\signame{busy} is high; the buffer is in \sv{CTRL\_COMPUTE} with
\signame{buf\_reg\_add}~=~0 and \signame{buf\_bit\_sel}~=~\signame{bit\_count}.
The state self-loops while \signame{pulse\_done} is low and advances to
\stateval{S\_RESULT} when \signame{pulse\_done} is high. \signame{bit\_count}
is incremented each cycle in \stateval{S\_COMPUTE} and cleared on exit.

\paragraph{\stateval{S\_RESULT}.}
\signame{busy} is high; the buffer is in \sv{CTRL\_RESULT\_OUT}. The
transition is unconditional: \signame{next\_state} becomes
\stateval{S\_IDLE}.

\paragraph{Default arm.}
On an illegal or unknown \signame{state} value, the RTL drives
\signame{next\_state}~=~\stateval{S\_IDLE} and
\signame{next\_prog\_state}~=~\sv{PROG\_HIZ}.

\subsection{Program Sub-FSM}
\label{sec:rtl-controller-prog}

The PROG sub-FSM has five states --- \sv{PROG\_HIZ}, \sv{PROG\_SELECT},
\sv{PROG\_WRITE}, \sv{PROG\_WAIT\_ACK}, \sv{PROG\_COMPLETE} --- and is
reset to \sv{PROG\_HIZ} on \signame{!rst\_n}.

The forward flow is summarized in \Cref{fig:prog-subfsm-body} and in
\Cref{tab:prog-subfsm-table}.

\begin{figure}[htbp]
    \centering
    \includegraphics[width=0.85\linewidth]{fig_prog_subfsm_body}
    \caption[Program sub-FSM]{Program sub-FSM inside \stateval{S\_PROGRAM}.
        Entry is from \stateval{S\_IDLE} via \sv{CMD\_PROG}; exit on
        \sv{PROG\_COMPLETE} sets \signame{next\_state}~=~\stateval{S\_VERIFY}.}
    \label{fig:prog-subfsm-body}
\end{figure}

\begin{table}[htbp]
    \centering
    \caption{Next-state table of the PROG sub-FSM.}
    \label{tab:prog-subfsm-table}
    \begin{tabularx}{\linewidth}{l X}
        \toprule
        \textbf{\signame{prog\_state}} & \textbf{Rule for \signame{next\_prog\_state}} \\
        \midrule
        \sv{PROG\_HIZ} & Unconditional: \sv{PROG\_SELECT}. \\
        \sv{PROG\_SELECT} & \sv{op\_done ? PROG\_WRITE : PROG\_SELECT}. \\
        \sv{PROG\_WRITE} & \sv{pulse\_done ? PROG\_WAIT\_ACK : PROG\_WRITE}. \\
        \sv{PROG\_WAIT\_ACK} & \sv{op\_done ? PROG\_COMPLETE : PROG\_WAIT\_ACK}. \\
        \sv{PROG\_COMPLETE} & Sets \signame{next\_state}~=~\stateval{S\_VERIFY}; sets \signame{next\_prog\_state}~=~\sv{PROG\_HIZ}. \\
        \sv{default} & \sv{PROG\_HIZ}. \\
        \bottomrule
    \end{tabularx}
\end{table}

\signame{op\_done} is a level-sensitive combinational condition: the RTL
contains no \sv{always\_ff} that samples \signame{op\_done} inside
\modname{ann\_controller}. The registered \signame{prog\_state} updates
on the next clock edge to whatever \signame{next\_prog\_state} resolves
to, which in turn depends on the current \signame{op\_done} value in the
same evaluation.

The \signame{pulses} bus is set to \sv{PULSE\_MODE\_PROG} only when
(\signame{state}~=~\stateval{S\_PROGRAM})~\sv{\&\&}~(\signame{prog\_state}~=~\sv{PROG\_WRITE})
and the corresponding train-active Boolean is true;
\signame{pulse\_cnt} increments only while the FSM stays in
\sv{PROG\_WRITE} and \signame{pulse\_done} is still low.

\subsection{Verify Sub-FSM}
\label{sec:rtl-controller-verify}

The VERIFY sub-FSM has four states: \sv{VERIFY\_IDLE}, \sv{VERIFY\_WAIT},
\sv{VERIFY\_CHECK}, \sv{VERIFY\_DONE}. The registered
\signame{verify\_state} is assigned each cycle from the combinational
\signame{verify\_next}. Outside of \stateval{S\_VERIFY}, the combinational
case does not assign \signame{verify\_next}; the shared \sv{always\_ff}
nevertheless samples it every cycle, which (in a fully-elaborated
\sv{always\_comb}) typically infers transparent latching of the previous
value when the main state is not \stateval{S\_VERIFY}.

The \sv{VERIFY\_IDLE} branch tests
\sv{PULSE\_TOTAL\_READ}~\(\le\)~\sv{8'd1}: if true, it short-circuits
directly to \sv{VERIFY\_CHECK}; otherwise, it enters \sv{VERIFY\_WAIT}.
With default constants in \sv{controller\_pkg}, \sv{PULSE\_TOTAL\_READ}
is 2, so the RTL takes the \sv{VERIFY\_WAIT} branch. The
\sv{VERIFY\_WAIT} branch compares the registered
\signame{verify\_pulse\_idx} against \sv{PULSE\_TOTAL\_READ[7:0]~-~8'd1}
and transitions to \sv{VERIFY\_CHECK} when they match; otherwise it
self-loops. The index \signame{verify\_pulse\_idx} is cleared to
\sv{8'd0} when the main state is not \stateval{S\_VERIFY}, loaded to
\sv{8'd1} in \sv{VERIFY\_IDLE}, and incremented by 1 in
\sv{VERIFY\_WAIT} when the next sub-state is also \sv{VERIFY\_WAIT}.

The \sv{VERIFY\_CHECK} arm is the comparator body and is described
separately in \Cref{sec:rtl-controller-verify-check}. On \sv{VERIFY\_DONE}
the combinational logic sets \signame{next\_state}~=~\stateval{S\_IDLE}
and \signame{verify\_next}~=~\sv{VERIFY\_IDLE}.

The simplified verify sub-FSM is shown in \Cref{fig:verify-subfsm-body}.

\begin{figure}[htbp]
    \centering
    \includegraphics[width=0.9\linewidth]{fig_verify_subfsm_body}
    \caption[Verify sub-FSM]{Verify sub-FSM inside \stateval{S\_VERIFY}.
        \sv{VERIFY\_CHECK} has three outcomes; see
        \Cref{sec:rtl-controller-verify-check}.}
    \label{fig:verify-subfsm-body}
\end{figure}

\subsubsection{\texttt{VERIFY\_CHECK}: Compare, Retry, and Erase Branch}
\label{sec:rtl-controller-verify-check}

In \sv{VERIFY\_CHECK}, the RTL compares the 4-bit
\signame{weight\_read\_data} input with \signame{expected\_weight}. The
package definition is \sv{assign expected\_weight = weight\_from\_buffer},
and within the combinational region \signame{weight\_from\_buffer} is
\signame{buf\_data}[3:0] when the state is \stateval{S\_PROGRAM} or
\stateval{S\_VERIFY}. Since \signame{buf\_addr\_reg} for those states is
\signame{address\_reg}[5:0], the expected nibble is the buffer contents
at the programmed cell address.

The three outcomes are:

\begin{enumerate}
    \item \textbf{Match:}
          \signame{weight\_read\_data}~=~\signame{expected\_weight} sets
          \signame{verify\_next}~=~\sv{VERIFY\_DONE}.
    \item \textbf{Under-programmed}
          (\signame{weight\_read\_data}~\(<\)~\signame{expected\_weight}):
          \begin{itemize}
              \item If \signame{prog\_retry\_cnt}~\(<\)~\sv{MAX\_PROG\_RETRIES}
                    (the package sets \sv{MAX\_PROG\_RETRIES}~=~3):
                    \signame{next\_state}~=~\stateval{S\_PROGRAM},
                    \signame{next\_prog\_state}~=~\sv{PROG\_HIZ}.
              \item Else:
                    \signame{verify\_failure\_starts\_erase}~=~1 and
                    \signame{next\_state}~=~\stateval{S\_ERASE}.
          \end{itemize}
          In both under-programmed sub-cases, \signame{verify\_next}
          becomes \sv{VERIFY\_IDLE}.
    \item \textbf{Over-programmed}
          (\signame{weight\_read\_data}~\(>\)~\signame{expected\_weight}):
          \signame{verify\_failure\_starts\_erase}~=~1,
          \signame{next\_state}~=~\stateval{S\_ERASE},
          \signame{verify\_next}~=~\sv{VERIFY\_IDLE}.
\end{enumerate}

The registered \signame{prog\_retry\_cnt} is updated in a priority
if/else-if chain inside the sequential block. When the main state is
\stateval{S\_VERIFY} and the sub-state is \sv{VERIFY\_DONE},
\signame{retry\_cnt}, \signame{prog\_retry\_cnt}, and
\signame{program\_stronger} are all cleared to 0. When the main state is
\stateval{S\_VERIFY}, the sub-state is \sv{VERIFY\_CHECK}, the readback
is less than expected, and \signame{prog\_retry\_cnt} is still below
\sv{MAX\_PROG\_RETRIES}, \signame{prog\_retry\_cnt} is incremented. When
the readback is above expected, or the retry budget is exhausted,
\signame{prog\_retry\_cnt} is cleared to 0. When the main state is
\stateval{S\_ERASE} and the erase sub-state is \sv{ERASE\_COMPLETE},
\signame{retry\_cnt} is incremented and \signame{prog\_retry\_cnt} is
cleared. A separate condition in the same \sv{always\_ff} asserts
\signame{program\_stronger} when \sv{VERIFY\_CHECK} shows
\signame{weight\_read\_data}~\(\ne\)~\signame{expected\_weight} and
\signame{expected\_weight}~\(>\)~\signame{weight\_read\_data}.

\signame{weight\_read\_en} is driven to 1 in \sv{VERIFY\_IDLE} and 0
elsewhere, but no other statement inside the module reads it; it is a
dead internal flag in the current RTL.

\subsection{Erase Sub-FSM}
\label{sec:rtl-controller-erase}

The ERASE sub-FSM has five states: \sv{ERASE\_HIZ}, \sv{ERASE\_SELECT},
\sv{ERASE\_PULSE}, \sv{ERASE\_WAIT\_ACK}, \sv{ERASE\_COMPLETE}. Its
forward flow is \sv{ERASE\_HIZ}\(\rightarrow\)\sv{ERASE\_SELECT}
(unconditional) \(\rightarrow\)\sv{ERASE\_PULSE} (requires
\signame{op\_done}) \(\rightarrow\)\sv{ERASE\_WAIT\_ACK} (requires
\signame{pulse\_done}) \(\rightarrow\)\sv{ERASE\_COMPLETE} (requires
\signame{op\_done}).

The exit branch from \sv{ERASE\_COMPLETE} is decided combinationally
from \signame{erase\_from\_host} and \signame{retry\_cnt}, as shown in
\Cref{tab:erase-complete}.

\begin{table}[htbp]
    \centering
    \caption{\texttt{ERASE\_COMPLETE} exit branches.}
    \label{tab:erase-complete}
    \begin{tabularx}{\linewidth}{l l l X}
        \toprule
        \textbf{Condition} & \textbf{\signame{next\_state}} & \textbf{\signame{next\_prog\_state}} & \textbf{Side effect} \\
        \midrule
        \signame{erase\_from\_host} & \stateval{S\_IDLE} & \sv{PROG\_HIZ} & --- \\
        \(!\)\signame{erase\_from\_host}, \signame{retry\_cnt}\,\(<\)\,3 &
            \stateval{S\_PROGRAM} & \sv{PROG\_HIZ} & --- \\
        Otherwise & \stateval{S\_IDLE} & \sv{PROG\_HIZ} &
            \signame{erase\_max\_retries\_exhausted}~=~1 \\
        \bottomrule
    \end{tabularx}
\end{table}

The \signame{erase\_from\_host} register distinguishes a host-issued
\sv{CMD\_ERASE} from a verify-initiated erase. Its update is a priority
if/else-if chain:
\begin{itemize}
    \item Cleared to 0 when
          (\signame{state}~=~\stateval{S\_VERIFY})~\sv{\&\&}~(\signame{next\_state}~=~\stateval{S\_ERASE})
          --- the verify-driven erase path.
    \item Set to 1 when
          (\signame{state}~=~\stateval{S\_IDLE})~\sv{\&\&}~\signame{valid}~\sv{\&\&}~(\signame{cmd}~=~\sv{CMD\_ERASE})
          --- a host erase command.
    \item Cleared to 0 when
          (\signame{state}~=~\stateval{S\_ERASE})~\sv{\&\&}~(\signame{next\_state}~=~\stateval{S\_IDLE})
          --- completion of the erase sequence.
\end{itemize}

The simplified erase sub-FSM is shown in \Cref{fig:erase-subfsm-body}.

\begin{figure}[htbp]
    \centering
    \includegraphics[width=0.9\linewidth]{fig_erase_subfsm_body}
    \caption[Erase sub-FSM]{Erase sub-FSM inside \stateval{S\_ERASE}.
        Exit from \sv{ERASE\_COMPLETE} branches according to
        \Cref{tab:erase-complete}.}
    \label{fig:erase-subfsm-body}
\end{figure}

\subsection{Pulse Generation and LUT-Based Programming}
\label{sec:rtl-controller-pulses}

The controller drives three distinct artefacts that together implement
pulse timing: a combinational per-state \signame{pulse\_total}, a
combinational \signame{pulse\_done}, and a registered
\signame{pulse\_cnt}.

\paragraph{Combinational \signame{pulse\_total}.}
The package helper \sv{pulse\_train\_total(T, N, G)} returns the total
cycle count for \(N\) bursts of \(T\) active cycles interleaved with
\((N-1)\) gaps of \(G\) \signame{HIZ} cycles. The controller assigns
\signame{pulse\_total} as an 8-bit unsigned clip of per-state values:

\begin{itemize}
    \item \stateval{S\_READ}: \sv{PULSE\_TOTAL\_READ} =
          \sv{pulse\_train\_total(TREAD, PULSE\_NUM\_READ, PULSE\_GAP)}.
    \item \stateval{S\_PROGRAM} with \signame{prog\_state}~=~\sv{PROG\_WRITE}:
          three branches, as summarized in \Cref{tab:prog-pulse-branches}.
    \item \stateval{S\_ERASE} with \signame{erase\_state}~=~\sv{ERASE\_PULSE}:
          \sv{PULSE\_TOTAL\_ERASE}.
    \item \stateval{S\_COMPUTE}: \sv{PULSE\_TOTAL\_INF}, which the package
          defines as \(\max(8, \mathrm{TINF} \cdot N + \text{gaps})\).
\end{itemize}

\begin{table}[htbp]
    \centering
    \caption{\signame{pulse\_total} branches in \stateval{S\_PROGRAM}
        \(+\) \sv{PROG\_WRITE}.}
    \label{tab:prog-pulse-branches}
    \begin{tabularx}{\linewidth}{l X}
        \toprule
        \textbf{Condition} & \textbf{Value} \\
        \midrule
        \sv{USE\_WEIGHT\_PULSE\_LUT}~=~0 &
            \sv{pulse\_train\_total(TPROG, PULSE\_NUM\_PROG, PULSE\_GAP)}. \\
        LUT enabled,
            \signame{prog\_retry\_cnt}~\(>\)~0 &
            Short retry train with \sv{Tp}~=~1, \sv{Np}~=~1. \\
        LUT enabled, first attempt &
            \sv{Mmacro}~=~\sv{pulse\_train\_total(TPROG, PULSE\_NUM\_PROG, PULSE\_GAP)};
            \sv{Rlut}~=~\sv{weight\_pulse\_cycles\_lut[expected\_weight]}
            (minimum 1); total
            \sv{pulse\_lut\_macro\_repeat\_total(Rlut, Mmacro, PULSE\_GAP)}. \\
        \bottomrule
    \end{tabularx}
\end{table}

\paragraph{Combinational \signame{pulse\_done}.}
The condition is
\sv{(pulse\_total~>~0) \&\& (pulse\_cnt~>=~pulse\_total~-~1)}. When
\signame{pulse\_total} is 0 (outside of the active states above),
\signame{pulse\_done} is always 0.

\paragraph{Registered \signame{pulse\_cnt}.}
The 8-bit counter is cleared to 0 on transitions into \stateval{S\_READ},
into \sv{PROG\_WRITE}, into \sv{ERASE\_PULSE}, or into
\stateval{S\_COMPUTE}, and it is incremented on cycles where the FSM
stays in the same pulse-active phase with \signame{pulse\_done} still
low. The pulse duration is therefore implemented as a load-and-count
counter compared against \signame{pulse\_total}.

\paragraph{Pulse-mode output \signame{pulses}[2:0].}
A separate combinational block drives \signame{pulses}. The default is
\sv{PULSE\_MODE\_HIZ} (\sv{3'b000}). For each active mode a Boolean
\emph{\_on} is formed from \sv{pulse\_train\_active} (for simple
burst/gap patterns) or \sv{pulse\_lut\_macro\_repeat\_active} (for
first-attempt LUT PROG). When the corresponding Boolean is true,
\signame{pulses} takes \sv{PULSE\_MODE\_READ}, \sv{PULSE\_MODE\_PROG},
\sv{PULSE\_MODE\_ERASE}, or \sv{PULSE\_MODE\_INF} as appropriate;
otherwise it is \sv{PULSE\_MODE\_HIZ}. The physical waveform on
\signame{pulses} is thus a time-multiplexed burst train with
\sv{PULSE\_GAP}-length \signame{HIZ} windows between bursts, and
\Cref{fig:pulse-train-model} depicts one such train for the standard
\sv{pulse\_train\_total}\,/\,\sv{pulse\_train\_active} pair.

\begin{figure}[htbp]
    \centering
    \includegraphics[width=0.9\linewidth]{fig_pulse_train_model}
    \caption[Standard pulse train model]{Pulse train parameterized by
        $T$, $N$, $G$. Active phases are wide; \sv{PULSE\_GAP} (\(G\))
        cycles separate adjacent bursts. The total length
        \sv{pulse\_train\_total}$(T,N,G)$ is matched against
        \signame{pulse\_cnt} to produce \signame{pulse\_done}.}
    \label{fig:pulse-train-model}
\end{figure}

\paragraph{LUT path (first PROG attempt).}
When \sv{USE\_WEIGHT\_PULSE\_LUT} is set and
\signame{prog\_retry\_cnt} is zero, the package function
\sv{pulse\_lut\_macro\_repeat\_total(Rlut, Mmacro, G)} returns the total
number of cycles occupied by \sv{Rlut} copies of a macro of length
\sv{Mmacro}, separated by \sv{G}-cycle gaps. The repeat count \sv{Rlut}
is the entry of \sv{weight\_pulse\_cycles\_lut[]} indexed by the
\emph{expected} weight nibble (with a minimum of 1 enforced), and the
array itself is loaded from \sv{WEIGHT\_PULSE\_LUT\_FILE} (the
\sv{target/Controller/programming\_inputs/weight\_pulse\_lut.mem} image
in this repository) via \sv{\$readmemh}. \Cref{fig:lut-macro-repeat}
illustrates this structure.

\begin{figure}[htbp]
    \centering
    \includegraphics[width=0.95\linewidth]{fig_lut_macro_repeat}
    \caption[LUT-based macro-repeat pulse structure]{LUT-based
        first-attempt PROG: the inner macro
        \sv{pulse\_train\_total(TPROG, PULSE\_NUM\_PROG, PULSE\_GAP)}
        is concatenated \sv{Rlut} times with \sv{PULSE\_GAP} cycles
        between copies, where \sv{Rlut} is
        \sv{weight\_pulse\_cycles\_lut[expected\_weight]} (minimum 1).}
    \label{fig:lut-macro-repeat}
\end{figure}

For \signame{prog\_retry\_cnt}~\(>\)~0, the RTL falls back to a short
\sv{Tp}~=~1, \sv{Np}~=~1 train, so subsequent programming attempts apply
single-cycle nudges rather than repeating the full first-attempt train.

\subsection{Inference Path}
\label{sec:rtl-controller-inf}

The inference path follows
\stateval{S\_COLLECT\_DATA}\(\rightarrow\)\stateval{S\_COMPUTE}\(\rightarrow\)\stateval{S\_RESULT}\(\rightarrow\)\stateval{S\_IDLE},
as shown in \Cref{fig:inference-dataflow}.

\begin{figure}[htbp]
    \centering
    \includegraphics[width=0.95\linewidth]{fig_inference_dataflow}
    \caption[Inference dataflow]{Inference dataflow. In
        \stateval{S\_COLLECT\_DATA}, eight pixel bytes are written into
        the input buffer in row-major order at
        \signame{row\_count}\(\cdot\)8\(+\)\signame{data\_count}. In
        \stateval{S\_COMPUTE}, the buffer exposes bit lanes
        \signame{D0}\ldots\signame{D7} using \signame{bit\_sel} stepped
        by the controller's \signame{bit\_count} from 0 to 7. Finally,
        \stateval{S\_RESULT} drives the buffer into \sv{CTRL\_RESULT\_OUT}
        for one cycle before returning to \stateval{S\_IDLE}.}
    \label{fig:inference-dataflow}
\end{figure}

\signame{op\_done} is not referenced anywhere on the inference path in
the current RTL: the handshake exists only in the PROG and ERASE
sub-FSMs. The duration of \stateval{S\_COMPUTE} is determined solely by
\signame{pulse\_done} relative to \sv{PULSE\_TOTAL\_INF}.

\subsection{Recovery Loop Summary}
\label{sec:rtl-controller-recovery}

\Cref{fig:recovery-loop-conceptual} collects the main-FSM transitions
that together form the closed-loop recovery policy required by the
write--verify problem. PROG enters VERIFY; VERIFY either succeeds and
returns to IDLE, re-enters PROG on an under-programmed readback within
the retry budget, or enters ERASE on an over-programmed readback or on a
retry-budget overflow. ERASE either returns to IDLE (host erase or
retry-count overflow) or re-enters PROG to retry.

\begin{figure}[htbp]
    \centering
    \includegraphics[width=0.95\linewidth]{fig_recovery_loop_conceptual}
    \caption[Program--Verify--Erase recovery loop]{Conceptual view of
        the closed-loop recovery policy implemented by the main FSM and
        its three sub-FSMs. The labels on the edges summarize the
        compound conditions described in
        \Cref{sec:rtl-controller-mainfsm,sec:rtl-controller-verify-check,sec:rtl-controller-erase}.}
    \label{fig:recovery-loop-conceptual}
\end{figure}

\subsection{Signals Assigned but Not Consumed}
\label{sec:rtl-controller-dead}

Three internal signals are assigned by the combinational block but are
not read anywhere else in \sv{controller.sv}:
\signame{verify\_failure\_starts\_erase} (set in \sv{VERIFY\_CHECK}
erase branches), \signame{erase\_max\_retries\_exhausted} (set on the
\sv{ERASE\_COMPLETE} give-up branch), and the already-noted
\signame{weight\_read\_en}. They may exist to support hierarchical
testbench access or future integration; in the shipped RTL they are
inert. The implications of these dead drivers are revisited in
\Cref{ch:limitations}.

\section{Input Buffer}
\label{sec:rtl-buffer}

The input buffer is described in
\sv{source/Input\_Buffer/input\_buffer.sv} and its package. It uses the
package constants \sv{BUFFER\_SIZE}~=~64 and \sv{BUFFER\_DATA\_WIDTH}~=~8,
giving a 64-entry byte store.

\subsection{Storage Array and Write Path}
\label{sec:rtl-buffer-write}

The storage is declared as
\sv{logic [BUFFER\_DATA\_WIDTH-1:0] buffer\_reg[0:BUFFER\_SIZE-1]}, a
vector of registers indexed by the 6-bit \signame{buf\_reg\_add} port.
The write-enable qualifier is
\signame{write\_en}~=~\signame{buf\_read\_write}~\sv{\&\&}~(\signame{reg\_ctrl}~=~\sv{CTRL\_DATA\_LOAD}).
On the positive edge of the clock, if \signame{write\_en} is asserted
and the address passes \sv{is\_valid\_addr}, then
\sv{buffer\_reg[addr]}~\(\leftarrow\)~\signame{data\_in}. An asynchronous
active-low reset initializes every entry to zero in a dedicated reset
loop.

\subsection{Byte Read Path (\signame{buf\_data})}
\label{sec:rtl-buffer-read-byte}

The read-enable qualifier is asserted when \signame{buf\_read\_write} is
low and \signame{reg\_ctrl} is one of \sv{CTRL\_COMPUTE},
\sv{CTRL\_RESULT\_OUT}, or \sv{CTRL\_WEIGHT\_READ}. A combinational
block drives \signame{buf\_data} to zero by default; if the read enable
is high and \sv{is\_valid\_addr(addr)} is true, the output becomes
\sv{buffer\_reg[addr]}. The byte-read path is therefore purely
combinational: a mux over the 64 storage entries selected by
\signame{buf\_reg\_add}.

\subsection{Bit-Serial Output Lanes \signame{D0}--\signame{D7}}
\label{sec:rtl-buffer-bitserial}

For each lane \(k \in \{0,\ldots,7\}\), the combinational block assigns

\[
    \signame{D}_k =
    \begin{cases}
        \sv{buffer\_reg[addr + k][bit\_sel]} &
        \text{if \signame{read\_en} and \(\text{addr}+k < \sv{BUFFER\_SIZE}\)} \\
        0 & \text{otherwise}
    \end{cases}
\]

Physically, \signame{buf\_reg\_add} is the \emph{base row} of an 8-pixel
group in the 1D array; the eight lanes present byte
\(\text{addr}{+}k\) to lane \(k\). The bit position inside each byte is
chosen by \signame{bit\_sel}[2:0], which the controller advances from 0
to 7 during \stateval{S\_COMPUTE} to implement LSB-first bit-serial
readout across the eight lanes.

\subsection{\signame{ready} Output}
\label{sec:rtl-buffer-ready}

\signame{ready} defaults to 0 and is asserted as \sv{is\_valid\_addr(addr)}
in exactly two mutually exclusive situations: (i)~when
\signame{reg\_ctrl}~=~\sv{CTRL\_DATA\_LOAD} and \signame{buf\_read\_write}
is 1 (write-load), and (ii)~when \signame{reg\_ctrl} is one of
\sv{CTRL\_COMPUTE}, \sv{CTRL\_RESULT\_OUT}, or \sv{CTRL\_WEIGHT\_READ}
and \signame{buf\_read\_write} is 0 (explicit reads). In every other
mode --- including \sv{CTRL\_IDLE} and any reads during
\sv{CTRL\_DATA\_LOAD} --- \signame{ready} is 0.

\subsection{Port Summary}
\label{sec:rtl-buffer-ports}

The complete port list of the input buffer is
\{\signame{clk}, \signame{rst\_n}, \signame{data\_in}, \signame{ready},
\signame{reg\_ctrl}, \signame{buf\_read\_write}, \signame{buf\_reg\_add},
\signame{bit\_sel}, \signame{buf\_data},
\signame{D0},\ldots,\signame{D7}\}.
Full port tables for all three modules, including directions and widths,
are collected in \Cref{app:port-tables}.
